\section{Discussion and Research Agenda}

\subsection{Theoretical Implications}

OAL represents trust as a structured relying relation. This yields three design implications. First, assurance arguments can be composed through typed dependencies when their inference rules and shared assumptions are explicit. Entity evidence supports authorization claims, device attestation supports claims about measured state, and observation and mission evidence support downstream use. Formal composition requires a theorem with explicit rules and assumptions for the combined claims. Second, trust is purpose-bound. Evidence sufficient for sharing low-sensitivity imagery may be insufficient for cooperative navigation, even when the underlying object is identical. Third, uncertainty requires an operational response: ``unknown'' records insufficient support for a conclusion and maps to a documented policy action. This epistemic status is represented separately from physical compromise and degradation.

The assurance chain also clarifies the relationship between provenance and truth. Provenance is an argument about history: who generated an artifact, through which activities, from which inputs. Observation validity is an argument about the world. Ledger consistency is an argument about shared records. Mission suitability is an argument about consequence. These claims require separate tests even when they share evidence or failure causes. Provenance and decision-accountability research supports traceable transformations \citep{singh2019decision,kale2023provenance}; OAL makes explicit the mission consumer's judgment about which history supports which use.

\subsection{System Resilience under Partition and Degradation}

Future architectures should treat partition as a routine operating mode. Local policy envelopes require bounded capability, expiry, evidence-age limits, safe fallback, and explicit records of which global facts were unavailable. Reconnection combines synchronization with semantic reconciliation: the system determines which policy was locally available at event time, when a revocation became effective, which dependent claims require reassessment, and whether conflicting observations require mission review. Authorization under a previously delegated envelope and compliance with subsequently learned global state are distinct questions. Blockchain scalability reviews examine throughput and storage among other constraints \citep{hafid2020scaling,khan2021systematic}; cross-domain missions also require explicit treatment of long disconnection, sparse contacts, and consequence-sensitive local autonomy.

Resilience also requires functional diversity. If the USV is simultaneously acoustic gateway, time source, fusion node, policy-enforcement point, and validator, redundant vehicle count can conceal a shared failure dependency. The graph identifies vertices that support many mission paths; quantifying failure probability or mission loss additionally requires reliability and consequence models. Future work should compare architectural redundancy (multiple instances) with evidential independence (different roots, media, sensors, and organizations). Two reports relayed through one gateway share transport and availability dependencies. Independently signed source observations may retain distinct origin evidence; whether the gateway can also alter or bias their interpretation depends on end-to-end protection, selection, and transformation roles.

\subsection{Governance, Privacy, and Human Authority}

Consortium governance is part of the security mechanism. Research should specify who can admit validators, approve reference values, activate contracts and models, invoke emergency powers, decide disputes, and audit the auditors. Maritime blockchain adoption studies already identify policy and organizational factors beyond technical performance \citep{yang2019maritime,zhou2020key,farah2024survey}. In mission assurance, a governance failure can propagate like a cyberattack: an unsafe reference value makes compromised devices appear healthy, or a delayed revocation keeps stale authority active.

Privacy questions extend beyond encrypting payloads. Cross-mission identifiers, validator endorsements, timing, route custody, and model-use records can reveal organizational relationships or vehicle activity. Work on blockchain and IoT warns of immutable metadata and access-control tension \citep{ali2019applications,reyna2018blockchain,kshetri2017can}. Future designs should evaluate unlinkability, data minimization, selective disclosure, authorization-policy confidentiality, and deletion semantics alongside audit value. A mission may need accountable pseudonymity: an operator can verify that an authorized role acted while the long-term identity remains protected, and a governed process can resolve accountability after an incident.

Human authority must remain explicit. Automated trust inference can recommend reduced privilege, corroboration, or isolation, while high-consequence recovery and dispute decisions may require human review. The relevant evidence comprises model output, explanation, policy version, uncertainty, alternatives, and consequence. Authorized and auditable human override preserves this evidence trail and the accountability of the relying decision.

\subsection{Testable Research Propositions}

The following hypotheses translate the synthesis into comparisons for future studies. Each test requires a declared mission distribution, threat model, implementation, primary outcome, and minimum practically meaningful effect. Counterexamples and null effects delimit the conditions under which the proposed structure is useful.

\textbf{P1 (assurance separation).} In matched scenarios containing correctly signed but physically invalid observations, an implementation retaining separate observation, provenance, ledger, and mission-suitability assessments will have a lower unsafe-acceptance rate than a single-score implementation at a matched false-isolation cost. Both receive the same raw evidence, action set, and tuning budget. The outcome is the fraction of invalid products accepted for a specified consequential use, evaluated against an independent physical reference.

\textbf{P2 (dependency awareness).} For a specified binary event and prediction horizon, a probabilistic implementation that models shared sensor and recommender causes will have lower calibration error under correlated error and collusion than its conditional-independence ablation. Both use the same evidence and separately fitted parameters. Evaluation uses a prespecified estimator of held-out calibration, a proper scoring rule such as the Brier score, and decision loss, with results reported across correlation strengths and levels of lineage completeness. The test applies to outputs with an explicit probabilistic meaning.

\textbf{P3 (diagnosis and uncertainty).} Under mixed nominal operation, environmental degradation, adversarial intervention, and evidence outages, an implementation retaining competing causal hypotheses and an explicit unknown outcome will achieve lower prespecified mission loss than a binary trusted/untrusted implementation. Both have access to the same evidence and operational actions. Separate ablations of causal diagnosis and the ability to abstain distinguish their contributions. The loss function includes corroboration delay, false isolation, unsafe acceptance, and lost availability, allowing the effects of diagnosis and response policy to be assessed separately.

\textbf{P4 (bounded offline authority).} Across prespecified partition durations and revocation schedules, bounded offline authorization will complete more permitted tasks before their deadlines than authorization that requires current global approval, and will incur fewer actions conflicting with effective global policy than unrestricted offline autonomy. Both outcomes and resource costs are reported as a trade-off curve over envelope duration and scope. Event-time local authorization and conflicts discovered after reconnection are evaluated separately, using explicit revocation-effective times and clock uncertainty.

\textbf{P5 (selective ledger value).} Under administrator equivocation or record withholding within a declared fault and network model, a permissioned ledger with independent validators will provide more complete, mutually verifiable audit histories than a single-administrator signed log. The comparison measures the fraction of previously accepted test events recovered, detection of inconsistent views, and time to establish an agreed history, using an independent event record as ground truth. A witnessed or replicated signed-log baseline tests whether any gain requires the ledger architecture. Communication cost and deadline misses are measured separately; substantive governance disputes require their own resolution procedures alongside record agreement.

\textbf{P6 (evaluation coverage).} A system configuration selected using transaction and network metrics alone will have higher held-out mission loss than one selected from the same candidates using observation, endpoint, governance, and mission outcomes in addition to those metrics. The comparison uses equal evaluation budgets, fixed selection rules, and paired held-out scenarios containing both benign and adversarial failures. The measured difference tests whether the added evaluation dimensions improve selection for that mission distribution.

Each comparison requires prespecified operating thresholds and analysis rules, effect sizes with uncertainty intervals, and reporting of null or adverse outcomes. Findings apply to the tested implementation and mission distribution; broader applicability requires evaluation across additional configurations and settings.

\subsection{From OAL Cells to an Assurance Case}

OAL can structure an assurance case when each high-consequence mission use is expressed as a scoped claim. For example, ``fused product $d_{17}$ is suitable for route replanning by consumer $r$ until time $t$ under weather class $w$ and policy $p$.'' The argument links that claim to required evidence across the layers. Observation subclaims identify the sensors, calibration, geometry, environmental model, residuals, and corroboration that support physical plausibility. Commitment subclaims identify protected input versions, producers, transformations, custody, and time. Where shared ledger state is required, ledger subclaims state membership, endorsement, ordering, and the agreed policy version. Mission subclaims connect freshness and uncertainty to an explicit consequence and permitted action. Missing support is recorded as an evidence gap; counter-evidence or a challenged assumption is recorded as a potential defeater. Both may justify restricting action, but their epistemic meanings differ.

This structure records both supported scope and boundary conditions. A verifier may establish that a measured software component matches an approved reference, with camera alignment and operator intent assigned to separate subclaims. A provenance record may establish the instrumented derivation path and identify a legacy transformation as a gap in completeness. A ledger may establish endorsement of a model version by three authorized peers, while safety and accuracy are supported by separate validation evidence. Pairing each guarantee with its conditions gives downstream consumers the exact scope of justified reliance.

Assurance cases must also be time-aware. Evidence has an observation time, issue time, receipt time, and expiry condition; a dependency can change between them. If a USV attestation is fresh when it receives a UUV batch but stale when it transforms and endorses that batch, the later activity requires an explicit continuity assumption or renewed evidence. Clock uncertainty is represented where GNSS is unavailable or suspected. For disconnected operation, the case records the last globally known policy and the local authorization envelope, preserving the scope of agreement available at event time.

Counter-evidence is first-class. A new sensor residual, revocation, incident report, or conflicting custody event should identify which claims and descendants become disputed. This enables proportional response: quarantine a derived track while preserving independently supported observations from a platform, or reduce a gateway's endorsement role while retaining its emergency relay function. Recovery closes the argument through review of affected credentials, software, models, data descendants, and mission decisions, producing a scoped and evidence-backed return to service.

Future tools could generate these arguments from the typed dependency graph while preserving human judgment. Policy authors decide evidence sufficiency, independence, consequence, and acceptable uncertainty. The generated artifact therefore distinguishes machine-verified bindings (signature, hash, sequence), verifier appraisals (measured state against reference), statistical assessments (calibration or anomaly), and governance judgments (waiver, dispute, emergency authority). A reviewer can then compare the conclusion with the strongest supported subclaim.

Finally, evaluation should measure the assurance case itself. Evidence coverage can be assessed over an independently reviewed set of consequence-relevant dependencies, with critical gaps reported individually alongside aggregate coverage. Seeded counterexamples test sensitivity to invalid arguments; formal soundness requires a proof under the stated inference rules. Operator studies can measure the time and error in identifying affected products; governance exercises can measure unresolved conflicts and time to an authorized decision. These measures complement detection and ledger performance by testing whether the structure helps users reach appropriately scoped relying decisions.

\subsection{Toward a Cross-Domain Unmanned Trust Profile}

The proposed \CDUTP{} is an agenda for an interoperability profile: a small, versioned set of claims and evidence references for OAL boundaries shared by organizations. Its development agenda covers the schema, protocol bindings, and conformance tests. NIST Zero Trust Architecture supplies resource- and request-centered policy principles \citep{rose2020zero}; IETF RATS supplies attestation roles and semantics \citep{birkholz2023remote}; W3C PROV-DM and PROV-O supply entity--activity--agent provenance relations \citep{w3c2013provdm,w3c2013provo}. NIST's IoT device baseline informs endpoint capabilities \citep{fagan2020iot}. 3GPP specifications for UAS support and maritime services provide communications and service context \citep{threegpp22125,threegpp22119}. Permissioned-ledger governance and endorsement provide candidate mechanisms for multi-party state coordination \citep{androulaki2018hyperledger,hyperledger2026fabric}. Conformance and interoperability require version-specific mappings and tests across the referenced architectures and specifications.

Software update evidence can reuse secure-update architecture. The IETF SUIT architecture describes manifests, authorization, payload acquisition, installation, and device constraints \citep{ietf2020suit}; CDUTP could reference an update identifier and installation result while keeping the firmware payload off chain. Applying similar bindings to learned models would additionally require model-specific compatibility and validation evidence. These mappings would connect domain claims to the mission assurance graph while preserving the scope of the underlying protocol guarantees.

A profile instance should minimally state: object identifier and type; claim and assurance layer; issuer and relying role; evidence or provenance reference; observed, issued, and expiry times; mission context and permitted use; uncertainty or appraisal outcome; applicable policy, model, schema, and contract versions; privacy label and disclosure conditions; parent dependencies; revocation reference; and on/off-chain binding. It should support disconnected issuance and later verification while labeling freshness at the relying time. Compact encodings are desirable for acoustic links, but semantic equivalence is more important than one wire format.

A staged standardization pathway begins with published use cases and claim-scope requirements, followed by a shared vocabulary and conformance tests for evidence lineage and version binding. Subsequent stages comprise plugfests across independently implemented UAV, USV, UUV, gateway, verifier, and ledger components; multi-organization privacy and governance evaluation, including member exit and disputed recovery; and alignment with relevant standards bodies. CDUTP provides a research agenda for that alignment.

\subsection{Design Heuristics}

Several design principles follow from the synthesis. Technology selection begins with the relying decision, object, and assurance claim, which define the roles of the trust method and shared-state service. Evidence management preserves evidence identity and common-cause information, stores raw evidence off chain, minimizes metadata, and explicitly represents staleness and unknown status. Architectures separate hard real-time safety from shared audit and treat policies, models, contracts, and reference values as versioned artifacts exposed to threats. Evaluation includes signed-log and replicated-database alternatives. Descriptions such as ``proposed,'' ``analyzed,'' ``simulated,'' ``prototyped,'' ``field-tested,'' and ``deployed'' identify distinct forms of support, with claim scope stated for each.

These heuristics guide technology selection. Signed append-only logs and replicated databases are candidate solutions when one accountable manager governs state and audit. Distributed agreement may be useful when independent parties require a common history, subject to the fault model, availability constraints, and comparison with independently witnessed logs. Evidence commitments require adequate key custody, time or sequence guarantees, and protected local state. Decisions with deadlines shorter than end-to-end communication require a local safety mechanism whose assumptions and fallback behavior are validated separately.

Responsible deployment adds one final boundary. Mission-trust infrastructure can improve safety, environmental observation, inspection, disaster response, and coalition accountability, while persistent identities and provenance traces also create surveillance and operational-security consequences. An assurance case should therefore include a legitimate-purpose claim, proportional data collection, retention and disclosure limits, and an accountable path for human appeal or incident review. Dual-use assessment focuses on capabilities and governance. Public reporting can disclose architecture, assurance semantics, failure modes, and evaluation methods while protecting credentials, exact operational configurations, exploitable topology details, and mission-specific tactics. This balance supports scientific testing and responsible dissemination. It also reinforces the core OAL insight: governance artifacts and privacy policies are versioned, contestable evidence within the mission system and require the same accountable appraisal as technical artifacts.
